\begin{abstract}
针对异构 AI 任务与新能源、储能及通信网络耦合问题，本文建立算—储—电协同调度模型。问题一采用固定起点季节基线与 Huber-Ridge 模型预测 GPU 需求，测试期预测总量为 13155.739 等效 GPU，MAE 均值为 25.837 等效 GPU；完成 538 个任务的末端调度，其中 68 个任务在收尾时域结清，表明预测与实际任务调度实现了分离衔接。问题二建立成本—碳排放—时延调度模型，以确定性可行启发式生成初始解，并对高 GPU-hour 弹性任务进行有界局部搜索，得到运行成本 -421711041 元、碳排放 188.761 tCO2、网络时延 5.348 ms及新能源利用率 67.96\textbackslash{}。问题三建立固定负荷下的储能能量流模型，采用确定性储能规则筛选可行候选，累计充、放电量分别为 2076.315 MWh和 0.0 MWh；成本增加 304224 元，绝对爬坡量增加 247.775 MW，说明该规则未改善经济性与波动性。问题四建立任务迁移、开工时段、储能及购售电联合多目标模型，采用确定性可行启发式与有界局部搜索形成 5 个权衡候选；推荐方案迁移 94 个任务，运行成本为 -421426260 元，碳排放为 0.238 tCO2，新能源利用率为 67.97\textbackslash{}。
\end{abstract}
